%=====================================================================
\chapter{Desktop and Application Virtualization}\label{ch:desktop}
%=====================================================================
\locator{4}{Part IV --- Storage, Network and Cloud Virtualization}

\begin{outcomes}
By the end of this chapter you will be able to:
\begin{tight}
  \item \textbf{Explain} what VDI moves and what stays, and \textbf{contrast}
        persistent with non-persistent desktops.
  \item \textbf{Describe} what a remote display protocol actually sends, and
        \textbf{explain} why latency matters more than bandwidth.
  \item \textbf{Compute} the bandwidth and storage a VDI deployment needs, and
        \textbf{identify} the resource that fails first.
  \item \textbf{Distinguish} desktop virtualization from application streaming
        and from remote browser isolation.
  \item \textbf{Decide} whether VDI is the right answer to a stated
        requirement, and \textbf{say} what it costs.
\end{tight}
\end{outcomes}

\begin{prereq}
\chref{storage} (thin provisioning and linked clones --- a non-persistent
desktop estate is the largest copy-on-write deployment most organisations ever
run) and \chref{cpumem} (overcommit, which VDI pushes harder than any server
workload).
\end{prereq}

\begin{keyterms}
VDI $\bullet$ virtual desktop $\bullet$ persistent desktop $\bullet$
non-persistent desktop $\bullet$ golden image $\bullet$ linked clone $\bullet$
instant clone $\bullet$ profile management $\bullet$ remote display protocol
$\bullet$ PCoIP $\bullet$ Blast $\bullet$ ICA/HDX $\bullet$ RDP $\bullet$
boot storm $\bullet$ IOPS $\bullet$ session host $\bullet$ published
application $\bullet$ application streaming $\bullet$ remote browser isolation
\end{keyterms}

\section{Moving the Computer, Leaving the Screen}

Everything so far has virtualized a server. \term{Virtual desktop
infrastructure} does it to the thing on somebody's desk: the operating system,
the applications and the files run in the data centre, and what travels to the
user is pixels in one direction and keystrokes and mouse movements in the
other.

The motives are the four of \chref{what}, pointed at a different problem:

\begin{tight}
  \item \textbf{Isolation and control} --- the data never leaves the data
        centre, which is the reason regulated industries buy it;
  \item \textbf{Consolidation} --- a hundred desktops on a few hosts rather
        than a hundred machines under a hundred desks;
  \item \textbf{Encapsulation} --- a broken desktop is replaced in ninety
        seconds, because it is a file;
  \item \textbf{Access} --- the same desktop from a laptop, a thin client or a
        tablet, from anywhere.
\end{tight}

\begin{definitionbox}[Persistent and Non-Persistent Desktops]
A \term{persistent} desktop is one virtual machine per user, kept between
sessions. Users install software, change settings and leave files on it,
exactly as on a physical PC. Familiar, flexible, and it must be patched,
backed up and managed one hundred times over --- the server sprawl of
\chref{what}, re-created virtually.

\smallskip
A \term{non-persistent} desktop is created from a \term{golden image} at logon
and destroyed at logoff. Every session starts identical and known. Patching
means updating one image. The user's identity has to be reassembled each time
from \term{profile management} --- documents on a file share, settings in a
profile container --- and anything not captured by that is lost at logoff,
which users discover the hard way.

\smallskip
Non-persistent is the design that makes the economics work, and profile
management is where it succeeds or fails. Almost every unhappy VDI deployment
is a non-persistent one whose profile handling was underestimated.
\end{definitionbox}

\begin{conceptbox}[Linked and Instant Clones: Chapter 12, Applied at Scale]
A thousand non-persistent desktops from one image are a thousand
copy-on-write overlays over one read-only base --- exactly the structure of
\chref{storage}, and of \chref{containers}'s overlayfs one layer up.

\smallskip
A \term{linked clone} shares the base disk and keeps a delta per desktop:
a 60 GB image plus perhaps 2 GB each, so a thousand desktops cost 2 TB rather
than 60 TB. An \term{instant clone} goes further and forks a \emph{running}
machine's memory as well, so a new desktop is available in one or two seconds
rather than ninety.

\smallskip
The read amplification warning of \chref{storage} applies in full. Every one of
those thousand desktops reads the same base through the same chain, so the base
must be on the fastest storage you have and must stay in cache. VDI is where
storage design stops being a background concern.
\end{conceptbox}

\section{What Travels Over the Wire}

A remote display protocol is not a video stream. It inspects what changed on
the screen and sends the cheapest representation of that change: a rectangle of
pixels, a command to move a region, a glyph, or --- when the content really is
video --- an H.264 or H.265 stream with the work done on a server GPU.

\dsfig{%
\begin{tikzpicture}[font=\scriptsize,
  bx/.style={rounded corners=2pt, line width=0.8pt, align=center,
             inner sep=3pt, minimum height=0.56cm, font=\scriptsize},
  dc/.style={bx, draw=primary, fill=primary!10, text=primary,
             minimum width=3.4cm},
  cl/.style={bx, draw=greenhl, fill=greenhl!10, text=greenhl!45!black,
             minimum width=3.0cm},
  note/.style={font=\scriptsize\itshape, align=center}
]
  \node[rounded corners=3pt, draw=primary!45, fill=primary!3,
        minimum width=4.4cm, minimum height=3.5cm] at (0,0) {};
  \node[font=\scriptsize\bfseries, text=primary] at (0,1.95) {Data centre};
  \node[dc] at (0,1.25) {the desktop guest};
  \node[dc] at (0,0.55) {its applications};
  \node[dc] at (0,-0.15) {its files};
  \node[dc, draw=goldhl, fill=goldhl!12, text=goldhl!45!black] at (0,-0.95)
       {encoder (often a GPU)};

  \node[rounded corners=3pt, draw=greenhl!45, fill=greenhl!3,
        minimum width=4.0cm, minimum height=3.5cm] at (10.6,0) {};
  \node[font=\scriptsize\bfseries, text=greenhl!45!black] at (10.6,1.95)
       {The user};
  \node[cl] at (10.6,1.25) {a screen};
  \node[cl] at (10.6,0.55) {a keyboard and mouse};
  \node[cl] at (10.6,-0.15) {a decoder};
  \node[cl, draw=neutral!55, fill=neutral!8, text=neutral!45!black, dashed]
       at (10.6,-0.95) {no data, ever};

  \draw[-{Stealth[length=2.4mm]}, goldhl, line width=1.2pt]
        (2.3,0.75) -- (8.5,0.75);
  \node[note, text=goldhl!45!black] at (5.4,1.15)
       {changed rectangles, glyphs, move commands,\\
        H.264 where the content is really video};
  \node[note, text=goldhl!45!black] at (5.4,0.35)
       {150 kb/s idle $\bullet$ 2 Mb/s office work $\bullet$ 10--30 Mb/s video};

  \draw[-{Stealth[length=2.4mm]}, secondary, line width=1.2pt]
        (8.5,-0.55) -- (2.3,-0.55);
  \node[note, text=secondary!80!black] at (5.4,-1.0)
       {keystrokes and mouse positions --- a few kb/s};

  \node[rounded corners=3pt, draw=accent!60, fill=accent!5,
        text=accent!70!black, font=\scriptsize, align=center, inner sep=4pt,
        text width=10.6cm] at (5.3,-2.4)
       {\textbf{Latency, not bandwidth, is what users feel.} Every keystroke
        makes a round trip before its character appears. Under about 50 ms
        this is invisible; by 150 ms typing feels detached; past 250 ms the
        desktop is described as broken no matter how much bandwidth is
        free.};
\end{tikzpicture}}%
{A remote desktop session. The asymmetry is extreme --- megabits down, kilobits
up --- and the quality the user perceives is governed by the small return
arrow, not the large one.}{vdi-path}

\begin{alertbox}[The two failure modes nobody provisions for]
\textbf{The boot storm.} Nine o'clock, and nine hundred people log in within
twenty minutes. Nine hundred desktops boot, read the same base image, load the
same profile and start the same antivirus scan. Storage that comfortably serves
the estate all afternoon delivers a fraction of what is suddenly demanded, and
the queue makes logons take fifteen minutes.

\smallskip
It is a storage \emph{IOPS} problem, not a capacity or a CPU problem, and it is
the single commonest reason a VDI pilot succeeds and the rollout fails: fifty
users never produce it and nine hundred do.

\smallskip
\textbf{The antivirus scan.} Every desktop scanning its disk at the same hour
multiplies the same effect, with no user to blame. Stagger it, or use an
agentless scanner that works from outside the guest.
\end{alertbox}

\begin{worked}[Sizing a 500-Seat Deployment]
Five hundred knowledge workers, 2 vCPU and 4 GB each, Windows desktops on an
image of 60 GB. Size the hosts, the storage and the network, and find what
fails first.

\smallskip
\textbf{Step 1 --- memory, which sizes the hosts.} With 10:1 vCPU
oversubscription acceptable for desktops --- they are idle most of the time ---
memory is again the constraint. Allow no memory overcommit for desktops, since
they are interactive and swapping is immediately felt:
\[ 500 \times 4 = 2{,}000 \text{ GB} \]
On hosts of 512 GB, reserving 15\% for the hypervisor and headroom:
\[ \left\lceil \frac{2{,}000}{512 \times 0.85} \right\rceil
   = \left\lceil \frac{2{,}000}{435} \right\rceil = 5 \text{ hosts} \]
Add one for N+1: \textbf{6 hosts}.

\smallskip
\textbf{Step 2 --- check CPU.} $500 \times 2 = 1{,}000$ vCPU on six hosts of 64
cores $= 384$ cores:
\[ \sigma_{\text{cpu}} = \frac{1{,}000}{384} = 2.6 \]
Comfortable for desktop workloads, which sit idle between keystrokes.

\smallskip
\textbf{Step 3 --- storage capacity, with linked clones.}
\[ \underbrace{60}_{\text{golden image}} + \underbrace{500 \times 2}_{\text{deltas}}
   = 1{,}060 \text{ GB} \]
Against 500 full copies at 60 GB, which would be 30 TB. A \textbf{28:1}
saving, and the reason non-persistent desktops are affordable.

\smallskip
\textbf{Step 4 --- storage IOPS, steady state.} A knowledge-worker desktop
averages about 10 IOPS, roughly 80\% writes:
\[ 500 \times 10 = 5{,}000 \text{ IOPS} \]
Any modern NVMe array handles this without noticing.

\smallskip
\textbf{Step 5 --- storage IOPS, at the boot storm.} A booting Windows desktop
demands around 100 IOPS for perhaps two minutes. Nine hundred -- here 500 --
logging in over 20 minutes means roughly 50 booting concurrently at any moment
if arrivals are even:
\[ 50 \times 100 = 5{,}000 \text{ IOPS extra} \]
But arrivals are \emph{not} even. If half the estate logs in during the busiest
five minutes:
\[ \frac{250 \text{ desktops} \times 2 \text{ min boot}}{5 \text{ min}}
   = 100 \text{ concurrent} \qquad 100 \times 100 = 10{,}000 \text{ IOPS} \]
plus the 5{,}000 steady state: \textbf{15{,}000 IOPS, three times the
all-day figure, for five minutes.}

\smallskip
\textbf{Step 6 --- network.} At 2 Mb/s per session for office work:
\[ 500 \times 2 = 1{,}000 \text{ Mb/s} = 1 \text{ Gb/s} \]
If a tenth are watching video at 15 Mb/s:
\[ 450 \times 2 + 50 \times 15 = 900 + 750 = 1{,}650 \text{ Mb/s} \]
A 10 Gb/s data-centre uplink is ample; a site's \emph{internet} link may not be,
and remote users' home connections certainly vary.

\smallskip
\textbf{Step 7 --- what fails first.} Not memory, not CPU, not capacity, not
average IOPS. \textbf{Peak IOPS at logon}, by a factor of three over everything
else, for five minutes a day. Size the storage for step 5 and the rest follows;
size it for step 4 --- which is what the average suggests --- and the estate
works beautifully except at nine o'clock, which is when everybody forms their
opinion of it.
\end{worked}

\section{Virtualizing the Application Instead}

Delivering a whole desktop is often more than the problem requires.

\term{Published applications} (Citrix Virtual Apps, Azure Virtual Desktop in
multi-session mode) run many users' sessions on one shared \term{session host}
operating system and present individual application windows rather than a
desktop. Far denser than VDI --- dozens of users per host rather than a handful
--- because the operating system is shared. The cost is the one from
\chref{containers}: shared kernel, so one user's badly-behaved application
affects the others, and software that assumes it owns the machine misbehaves.

\term{Application streaming} goes further and delivers the application to the
user's own device, in a container that keeps it from touching the local
installation.

\term{Remote browser isolation} applies the idea to the one application
everybody runs. The browser executes in the data centre or at a provider, and
only rendered output reaches the endpoint, so a malicious page compromises a
container that is discarded afterwards. It is a security control rather than a
productivity one, and \chref{security} returns to it.

\rowcolors{2}{white}{lightbg}
\begin{longtable}{@{}L{3.4cm}L{3.5cm}L{3.5cm}L{4.0cm}@{}}
\toprule
\rowcolor{primary!15}
& \textbf{Isolation} & \textbf{Density} & \textbf{Use it when} \\
\midrule
\endhead
Persistent VDI & One guest per user & Lowest & Users need to install software,
or have genuinely individual environments \\
Non-persistent VDI & One guest per session & Moderate & A standard desktop for
many similar users; patching one image \\
Published apps & Shared session host & Highest & A few applications, many
users, no desktop needed \\
Application streaming & Container on the endpoint & N/A & Local performance
with central control of versions \\
Browser isolation & Disposable container per session & High & The threat is
the web, not the application \\
\bottomrule
\end{longtable}

\begin{conceptbox}[When VDI Is the Wrong Answer]
VDI is expensive, operationally demanding and unforgiving of under-sizing. It
is worth it when at least one of these is true:

\begin{tight}
  \item data must not leave the data centre, for regulatory or contractual
        reasons;
  \item users are numerous, similar, and need a standard environment;
  \item access is needed from many places or from untrusted devices;
  \item the endpoints are thin clients with a long life and little to steal.
\end{tight}

\smallskip
It is the wrong answer when users are few, when their needs are individual,
when the work is graphically heavy without a GPU budget (\chref{gpu}), or ---
most often --- when the real problem is patch management and a better endpoint
management tool would cost a tenth as much.

\smallskip
The question to ask of any VDI proposal is: \emph{what is it that the data
centre must hold?} If there is a crisp answer, VDI is likely right. If the
answer is ``it would be tidier'', it is not.
\end{conceptbox}

\begin{pitfall}
\begin{tight}
  \item \textbf{Sizing storage on average IOPS.} The boot storm is three times
        the average and it decides the user's opinion.
  \item \textbf{Piloting with fifty users and extrapolating.} Fifty users never
        produce a boot storm. Pilot the peak, not the mean.
  \item \textbf{Underestimating profile management.} Non-persistent desktops
        fail on this more than on anything technical.
  \item \textbf{Optimising bandwidth and ignoring latency.} Users feel the
        round trip on every keystroke. A fat, distant link is worse than a thin,
        near one.
  \item \textbf{Overcommitting desktop memory as if they were servers.} A
        swapping desktop is noticed instantly by the person using it.
  \item \textbf{Running in-guest antivirus on a schedule.} Every desktop
        scanning at once reproduces the boot storm with no users to show for
        it.
  \item \textbf{Treating VDI as a cost-saving measure.} It moves cost from
        endpoints to the data centre and adds licensing. Buy it for control,
        access or compliance, and be pleased if it also saves money.
\end{tight}
\end{pitfall}

\begin{lab}[A Remote Desktop, and What It Sends]
Builds a remote desktop session and measures what crosses the wire under three
kinds of work, then adds latency and sees which matters. Expect thirty-five
minutes.

\begin{lstlisting}[language=bash]
#!/usr/bin/env bash
# ch15_vdi.sh -- a remote desktop, measured, then made to feel slow.
set -euo pipefail

# --- 1. A desktop guest, from the Chapter 1 image. ---------------------
sudo virsh start vss1 || true ; sleep 25
ip_of () {
  sudo virsh domifaddr "$1" | awk '/ipv4/{split($4,a,"/"); print a[1]}'
}
IP=$(ip_of vss1)
G () { ssh -o StrictHostKeyChecking=no "ubuntu@$IP" "$1"; }
G "sudo apt-get -qq install -y xfce4 xrdp firefox && sudo systemctl \
   enable --now xrdp"

# --- 2. Connect with an RDP client and leave it idle. -----------------
# From the host: xfreerdp /v:$IP /u:ubuntu, or any RDP client.
sudo apt-get install -y freerdp2-x11 iftop tcpdump
echo "connect now, then press return"; read -r _

# --- 3. Measure the three cases of the figure. ------------------------
measure () {
  echo "=== $1 ==="
  sudo timeout 20 tcpdump -ni any "host $IP and port 3389" -q 2>/dev/null \
    | wc -l
}
measure "idle -- do nothing for 20 s"
echo "now type continuously in a text editor, then press return"; read -r _
measure "typing"
echo "now play a video full screen, then press return"; read -r _
measure "video"
# Packet counts stand in for bandwidth; the ratios are what matter.

# --- 4. Now the thing users actually feel: latency. -------------------
# Add 150 ms each way and use the session again. Nothing else changes.
sudo tc qdisc add dev virbr0 root netem delay 150ms
echo "use the session now -- type a sentence. Press return when done."
read -r _
sudo tc qdisc del dev virbr0 root
echo "and again with no added delay. The bandwidth was identical."

# --- 5. The boot storm, at small scale. -------------------------------
# Clone four desktops from one base and start them together; watch IOPS.
for n in 1 2 3 4; do
  sudo qemu-img create -f qcow2 -b /var/lib/libvirt/images/vss1.qcow2 \
       -F qcow2 "/var/lib/libvirt/images/d$n.qcow2" 20G
done
ls -lh /var/lib/libvirt/images/d*.qcow2   # linked clones: nearly empty
iostat -x 2 2 > /tmp/before.txt
for n in 1 2 3 4; do
  sudo virt-install --name "d$n" --memory 1024 --vcpus 1 \
    --disk "/var/lib/libvirt/images/d$n.qcow2",format=qcow2 \
    --import --graphics none --noautoconsole --os-variant ubuntu22.04 &
done
wait ; iostat -x 5 3 | tail -12      # the storm, four desktops wide

# --- 6. Tear down. ----------------------------------------------------
for n in 1 2 3 4; do
  sudo virsh destroy "d$n" 2>/dev/null || true
  sudo virsh undefine "d$n" --remove-all-storage 2>/dev/null || true
done
sudo virsh destroy vss1 || true
\end{lstlisting}

\textbf{What to notice.} Step 4 is the chapter. The bandwidth in that test was
exactly what it was in step 3, and the session became unusable anyway. Then
look at step 5's \texttt{iostat} output and multiply by a hundred.
\end{lab}

\begin{checkpoint}
\begin{tightnum}
  \item A 300-seat non-persistent deployment uses a 70 GB golden image with 3 GB
        deltas. Give the storage capacity needed, and the saving against full
        clones.
  \item At 10 IOPS steady state per desktop, and 120 IOPS for a 90-second boot,
        estimate peak IOPS if 150 of those 300 users log in within four
        minutes. Compare with the steady-state figure.
  \item A user reports that their remote desktop is ``laggy'' although the
        monitoring shows the link is 80\% idle. Explain what is probably
        happening and what you would measure.
\end{tightnum}
\end{checkpoint}

\begin{chaptersummary}
\begin{tight}
  \item VDI runs the desktop in the data centre and sends pixels out,
        keystrokes back. Buy it for control, access or compliance.
  \item Persistent desktops are familiar and must be managed individually;
        non-persistent ones are built from a golden image at logon and succeed
        or fail on profile management.
  \item Linked and instant clones are \chref{storage}'s copy-on-write at scale
        --- a 28:1 capacity saving in the worked example, and a read-amplified
        base that must stay fast.
  \item The display protocol sends changed regions, not video. Bandwidth is
        asymmetric and modest; \emph{latency} is what users feel, on every
        keystroke.
  \item Peak logon IOPS is three times the all-day figure and is what sizing
        must be based on. It is why pilots succeed and rollouts fail.
  \item Published applications are denser and share a kernel; browser isolation
        applies the idea as a security control.
\end{tight}
\end{chaptersummary}

\begin{reviewq}
  \item \textbf{(MCQ)} The resource that most often limits a VDI deployment
        is: (a)~host memory; (b)~peak storage IOPS at logon; (c)~network
        bandwidth; (d)~CPU.
  \item \textbf{(MCQ)} A non-persistent desktop keeps a user's documents and
        settings by: (a)~storing them in the golden image; (b)~profile
        management on separate storage; (c)~the delta disk; (d)~it does not
        keep them.
  \item \textbf{(Short)} Explain why latency matters more than bandwidth in a
        remote desktop session, in terms of what happens to a keystroke.
  \item \textbf{(Short)} Describe a boot storm and give two ways to reduce it.
  \item \textbf{(Short)} State the trade between published applications and
        full VDI, naming the chapter whose argument it repeats.
  \item \textbf{(Applied)} Size a 1{,}200-seat deployment of 2 vCPU, 8 GB
        desktops on 768 GB hosts with a 80 GB image and 4 GB deltas. Give host
        count, capacity, steady and peak IOPS, and bandwidth, and name what
        fails first.
  \item \textbf{(Applied)} A faculty proposes VDI so that students can use
        specialist software from their own laptops. Write the questions you
        would ask, the two alternatives you would price against it, and the
        circumstances under which you would recommend each.
\end{reviewq}
